\cotefile{136}
\scanpage{17}\mspage{15}
\deleted{\unsure{5)|g)}}\unsure{6)|h)} On peut associer à un espace topologique $X$ \deleted{\unsure{bien}} ou un ens $\frac12$simplicial $X_*$ \deleted{\unsure{des}}divers invariants alg. ``linéaires'' \uline{plus fins} à priori que le $C^{\cdot\cdot}_{\mathrm{DRpd}}$ (ou \unsure{$=$ que|mieux par} $C_{\cdot\cdot}^{\mathrm{DRpd}}$, \unsure{j'en dis autant|j'en dirais autant|en dualisant}\ldots).

[P. ex. on peut observer que $C^{\cdot\cdot}_{\mathrm{DRpd}\,*}(k)\simeq_d\Gamma^{\cdot}\Phi_{*k}\otimes_k\bigwedge^{\cdot}\Psi_{*k}$ se déduit de $D^{\cdot\cdot}_*(k)\unsure{=|/|\overset{\mathrm{déf}}{=}}\Gamma^{\cdot}\Phi_{*k}\otimes\bigwedge^{\cdot}\Phi_{*k}$ (qui est une Alg. diff \inserted{\unsure{homogène|$k$-augm.}} à p. div. qui est une \uline{résolution de} $k$) et de $T\in\Gamma(D_*^{10}{}_{k})$ comme Coker de la multiplication par $dT$ (i.e. on divise par l'idéal engendré par $dT$), ou encore en bidegré donné \deleted{\unsure{$p,q+1$}} $(p,q)$
\[
\deleted{\unsure{D}}C^{p,q}_{*k}\simeq\operatorname{Ker}\bigl(D_*^{p,q\deleted{\unsure{+1}}}\xrightarrow[\text{produit par }dT]{}D_*^{p,q+2}\bigr)
\]
d'où
\[
C^{p,q}_{\mathrm{DRpd}}(X_*,k)\simeq\operatorname{Ker}\bigl(D^{p,q+1}(X_*,k)\unsure{-|\longrightarrow}D^{p,q+2}(X_*,k)\bigr),
\]
et on peut considérer les structures \inserted{de base} sur $C^{\cdot\cdot}_{\mathrm{DRpd}}(X_*,k)$ comme \unsure{déduite|déduites} de certaines structures (à expliciter\ldots) sur $D^{\cdot\cdot}(X_*,k)$.

On peut aussi définir ``\uuline{La} structure multiplicative à p.d \inserted{cohomologique} du type d'homotopie $X_*$'', en \inserted{un sens convenable}, qui permet de reconstituer \deleted{\unsure{les}} \inserted{aussi bien} $C^{\cdot\cdot}_{\mathrm{DRpd}}(X_*,k)$ que $D^{\cdot\cdot}_k(X_*,k)$ (ou les complexes de De Rham-Sullivan). J'ignore si \unsure{le|la|les}
% archive sheet number: 16

